\begin{theorem}[Construct a good hierarchy or finish in the elementary branch]
Let \(G\) be a finite loopless multigraph with its individual edge type,
let \(e_0\) be the distinguished edge, and let the palette have \(k\)
colors.  Let
\[
 \mathcal E:\noderef{ExtendedTashkinovTree}(G,k,e_0,\rho,T,\varphi)
\]
be a positive-rung ETT with \noderef{MaximumProperty}, rung number \(n\), generating
coloring \(\varphi\), displayed tree \(T\), terminal ladder tree
\[
 T_n=\mathcal E.\mathit{series.data.tree}
       (\operatorname{Fin.last}(n).\mathit{castSucc}),
\]
and final series output
\[
 F=\mathcal E.\mathit{series.data.tree}
       (\operatorname{Fin.last}(n+1)).
\]
Assume a \noderef{TechnicalInductionContext} at rung \(n\),
\[
             \noderef{MaximumDegree}(G)+1\leq k,
             \qquad 11\leq \noderef{TashkinovOrder}(G).        \tag{1}
\]
Then one of the following alternatives holds.

First, \(V(T)\) is
\noderef{ColoredVertexSetElementary} under \(\varphi\).

Second, there are an ordered tree-sequence \(U\), a certificate
\[
 C_U:\noderef{CorrespondingExtendedTashkinovTree}
       (G,k,e_0,\mathcal E,\varphi,U)                            \tag{2}
\]
relative to the fixed original ETT \(\mathcal E\),
\noderef{MaximumProperty} for
\(C_U.\mathit{replacementETT}\), a generating-color
\noderef{HierarchyBase} \(B_U\) for that replacement ETT, a
\noderef{TashkinovHierarchy} \(H\) based on \(B_U\), and a
\noderef{GoodHierarchy} witness on \(H\), such that
\[
 \noderef{TreeColorSetClosed}
   \bigl(G,\varphi,U,\noderef{TreeMissingColors}(\varphi,U)\bigr) \tag{3}
\]
and
\[
                       V(U)=V(F).                               \tag{4}
\]
Equality (4) is only an equality of vertex sets.  No equality of the two
ordered tree-sequences, and no containment of an intermediate tree in
\(F\), is asserted.
\end{theorem}
\begin{proof}
Write \(a=\operatorname{Fin.last}(n)\), so that the exact terminal ladder
tree is \(T_n=\mathcal E.\mathit{series.data.tree}(a.\mathit{castSucc})\).
The terminal certificate supplied by \noderef{TashkinovSeries} is a
positive-index constructor of \noderef{TashkinovSeriesStep}.  In each of
its PE, SE, and RE cases, its current-coloring clauses, together with the
generating-color equality in \noderef{ExtendedTashkinovTree}, say that
\(\varphi\) is an \noderef{IsProperEdgeColoring} and is
\noderef{IsOneEdgeUncolored} at \(e_0\).  Reflexivity in \noderef{sc2}
makes \(\varphi\) a \noderef{StableColoring} of itself on \(T_n\), relative
to the terminal \(\noderef{TashkinovStoredColors}\).  The
\texttt{currentSix} field of the assumed
\noderef{TechnicalInductionContext}, instantiated with the original
arbitrary input \(\rho\), the equality \(\mathcal E.\mathit{rounds}=n\),
the maximum property, and these three coloring facts, therefore gives
\[
                       \noderef{ModColoring}(\mathcal E,\varphi). \tag{5}
\]
No mod-coloring is assumed.

Apply \noderef{HierarchyBaseExists} to \(\mathcal E\) and choose a
generating-color base \(B\).  Independently, apply
\noderef{TechnicalCurrentLadderMissingBoundLarge} to the maximum property,
the same technical context, and (1).  It gives the unconditional estimate
\[
                       2n+11\leq
       |\noderef{TreeMissingColors}(\varphi,T_n)|.               \tag{6}
\]
This invocation uses no hierarchy, hierarchy base, or prior missing-color
estimate.

We construct an initial state for a finite iteration, splitting only on the
actual terminal constructor.

Suppose first that its kind is PE.  Apply
\noderef{ParallelSpecialClosurePreparation} to \(\mathcal E\), the chosen
base \(B\), the maximum property, the technical context, (1), and (6).  Its
statement first supplies an exact \(\varphi\)-
\noderef{TashkinovAugmentationSequence}
\[
                         T_n\longrightarrow B.\mathit{core}      \tag{7}
\]
and the strict inequality
\(|T_n|<|B.\mathit{core}|\).  If that theorem returns elementarity of
\(V(T)\), we take the first alternative of the present theorem and stop.
This is exactly the early branch allowed before the paper imposes the
nonclosed-core hypothesis.

In the other PE alternative, the same theorem supplies a tree \(U_0\), a
certificate \(C_0\) already relative to the fixed original
\(\mathcal E\), \noderef{MaximumProperty} for its replacement ETT, a route
\[
                         B.\mathit{core}\longrightarrow U_0,     \tag{8}
\]
an exact initial-segment relation, and
\(|B.\mathit{core}|<|U_0|\).  It also supplies a base \(B_0\), a
zero-divider hierarchy \(H_0\), and a good witness
\(\mathcal G_0\).  The core, closure, and numerical support position of
\(B_0\) are explicitly equal to those of \(B\), and the exact two stages of
\(H_0\) are \(T_n\) and \(U_0\).  Finally it states that
\[
 U_0\text{ is }
 \left(\bigcup_{\eta\in
   \noderef{HierarchyTerminalStoredColors}(H_0)}
   \mathcal G_0.\Gamma^0_\eta\right)^-
 \text{-closed under }\varphi.                                 \tag{9}
\]
For this branch fix the anchor
\(J=B.\mathit{core}\).  Equations (7)--(9), including the strict inequality,
are the complete initial iteration state.  In particular, only the
zero-divider package in the second alternative is iterated.

Now suppose the terminal kind is not PE.  Positivity excludes the initial
constructor, so the terminal certificate in
\noderef{TashkinovSeriesStep} is SE or RE.  Apply
\noderef{AlgorithmFiveFiveBaseStage} with the maximum property, (5), the
base \(B\), and (6).  Each hypothesis whose antecedent says that the kind
is PE is discharged by the present non-PE assumption; no PE fact is used.
The theorem returns a seed \(J\), a tree \(U_0\), and a corresponding
certificate \(C_0\) relative to the fixed original \(\mathcal E\).  It
also returns \noderef{MaximumProperty}, a base \(B_0\) with the same core, closure,
and numerical support position as \(B\), a zero-divider hierarchy \(H_0\),
a good witness \(\mathcal G_0\), its terminal minus-closedness assertion,
and an exact route
\[
                              J\longrightarrow U_0.              \tag{10}
\]
Crucially, its non-PE seed clause retains one individual edge identity
\(f_n\) for which
\[
 \mathcal E.\mathit{series.data.connectingEdges}(a)=\{f_n\},
 \qquad \noderef{TreeSequenceAdjoinEdge}(T_n,f_n,J),             \tag{11}
\]
and it gives \(|T_n|<|U_0|\).  We keep this exact \(J\), rather than merely
its vertex set, fixed throughout the iteration.

We now describe the common finite iteration in either branch.  For a vertex
finset \(X\subseteq V(G)\), call \(X\) a state if there are
\[
 (U,C_U,B_U,H_U,\mathcal G_U)                                   \tag{12}
\]
such that \(V(U)=X\), \(C_U\) is corresponding relative to the fixed
original \(\mathcal E\), its replacement ETT has
\noderef{MaximumProperty}, the
three base data of \(B_U\) agree with those of \(B\), and \(H_U\) has good
witness \(\mathcal G_U\).  Require also the fixed route
\(J\longrightarrow U\), the terminal
\noderef{TreeColorSetMinusClosed} assertion
\[
 U\text{ is }
 \left(\bigcup_{\eta\in
   \noderef{HierarchyTerminalStoredColors}(H_U)}
   \mathcal G_U.\Gamma^{q_U}_\eta\right)^-
 \text{-closed},                                                \tag{13}
\]
where \(q_U=\operatorname{Fin.last}(H_U.\mathit{dividerCount})\),
and the fixed anchor alternative and strict inequality: in PE, (7),
\(J=B.\mathit{core}\), and
\(|B.\mathit{core}|<|U|\); in SE or RE, (11) and
\(|T_n|<|U|\).  The PE and non-PE constructions above show directly that
\(V(U_0)\) is a state.

Call \(X\) terminal when the full closedness condition (3) holds for any,
equivalently every, tree-sequence whose displayed vertex set is \(X\).
This is well defined: by the definitions of
\noderef{TreeSequenceVertexSet}, \noderef{TreeMissingColors}, and
\noderef{TreeColorSetClosed}, both the set of tree-missing colors and every
relevant colored boundary depend only on the displayed vertex finset.

Let \(X=V(U)\) be a nonterminal state, witnessed by (12).  The preceding
extensional observation says exactly that
\[
 \neg\noderef{TreeColorSetClosed}
       (G,\varphi,U,\noderef{TreeMissingColors}(\varphi,U)).      \tag{14}
\]
We verify the candidate-uniform bound required by
\noderef{AlgorithmFiveSixFixedRootRefresh}.  Given any arbitrary-input
maximum corresponding replacement ETT \(\mathcal L\) of rung \(n\), with
any base and good hierarchy, discard those last two objects and apply
\noderef{TechnicalCurrentLadderMissingBoundLarge} directly to
\(\mathcal L\).  The rung equality transports the fixed technical context
to the required type; the degree and large-order assumptions remain (1).
Thus the result is precisely
\[
            2n+11\leq
            |\noderef{TreeMissingColors}(\psi,L_n)|,             \tag{15}
\]
for the generating coloring \(\psi\) and exact terminal ladder tree
\(L_n\).  This proves the uniform premise without assuming a hierarchy in
the missing-color theorem.

Apply \noderef{AlgorithmFiveSixFixedRootRefresh} to the state witness,
(13)--(15), the fixed base \(B\), anchor \(J\), anchor alternative, strict
anchor inequality, and fixed route.  Its terminal output provides a tree
\(U^+\), a certificate \(C_{U^+}\) relative to the same fixed original
\(\mathcal E\), \noderef{MaximumProperty}, a transported base with the same three
base data, a hierarchy \(H_{U^+}\), a good witness, the corresponding
minus-closedness assertion, and a route
\(J\longrightarrow U^+\).  It also gives
\[
                              |U|<|U^+|.                         \tag{16}
\]
The nesting fields in the returned
\noderef{TashkinovHierarchy} say that \(U\) is an exact
\noderef{TreeSequenceInitialSegment} of \(U^+\).  Applying
\noderef{TreeSequenceInitialSegmentVertexSetStrict} to that nesting and
(16) gives
\[
                              V(U)\subsetneq V(U^+).              \tag{17}
\]
All other output fields are exactly the state conditions, with the same
fixed \(J\).  Hence every nonterminal state finset has a strictly larger
state finset.

Apply \noderef{FiniteStrictGrowthReachesTerminal} to this state predicate,
terminal predicate, strict-growth rule (17), and the initial state
\(V(U_0)\).  Since \(V(G)\) is finite, it returns a terminal state.  Choose
its complete witness (12) and call its tree \(U\).  We have obtained the
certificate (2), \noderef{MaximumProperty}, the base, hierarchy, good witness, and
full closedness (3).  It remains only to prove (4).

By \noderef{TashkinovClosureOfRouteAndAllMissingClosed}, the fixed route
\(J\longrightarrow U\) and (3) make \(U\) a
\(\varphi\)-\noderef{TashkinovClosure} of \(J\).  We compare it with the
original final output only now, after iteration has terminated.

In PE, concatenate the exact route (7), from \(T_n\) to the fixed anchor
\(J=B.\mathit{core}\), with the fixed-anchor route from \(J\) to \(U\).
The reflexive-transitive-chain definition of
\noderef{TashkinovAugmentationSequence} gives a route
\(T_n\longrightarrow U\), and (3) therefore makes \(U\) a closure of the
exact \(T_n\).  The PE constructor of the original terminal
\noderef{TashkinovSeriesStep} states that \(F\) is another
\(\varphi\)-closure of this same exact \(T_n\).  Apply
\noderef{TashkinovClosureVertexSetUnique} to these two closures to obtain
\(V(U)=V(F)\).

In SE or RE, the corresponding constructor of the original terminal
\noderef{TashkinovSeriesStep} supplies an exact edge \(g_n\), its own
one-edge seed \(J_0\), the equations
\[
 \mathcal E.\mathit{series.data.connectingEdges}(a)=\{g_n\},
 \qquad \noderef{TreeSequenceAdjoinEdge}(T_n,g_n,J_0),
\]
and a \(\varphi\)-closure of \(J_0\) equal to \(F\).  Comparison with
(11) gives equality of the singleton sets \(\{g_n\}=\{f_n\}\), hence
equality of the individual exact edges \(g_n=f_n\).  Rewrite the
constructor's adjoining certificate by this edge equality.  The theorem
\noderef{TreeSequenceAdjoinEdgeUnique} therefore gives \(J=J_0\) as full
ordered edge-valued tree-sequences.  Rewriting the original closure by this
equality makes \(F\) a closure of the fixed anchor \(J\).  The endpoint
tree \(U\) is also a closure of \(J\), so
\noderef{TashkinovClosureVertexSetUnique} again gives \(V(U)=V(F)\).

Thus (4) holds in every terminal-kind case.  Every exposed corresponding
certificate was returned relative to the fixed original \(\mathcal E\) by
\noderef{ParallelSpecialClosurePreparation},
\noderef{AlgorithmFiveFiveBaseStage}, or
\noderef{AlgorithmFiveSixFixedRootRefresh}; no nested replacement provider
is used in this proof.  Consequently no uncomposed intermediate provenance
remains, and all objects in the second alternative have the claimed types.
\end{proof}
